\scanpage{148}
\mspage{97}

trouver un exemple où elle est satisfaite, sans que $U_{/H}$ elle-même ne soit tot. $\unsure{W|\mathcal W}$-asphérique.

L'existence d'objets \deleted{tot.} totalement $W$-asphériques dans $\mathrm{Cat}$ n'est pas tellement évidente a priori. Qu'il y \inserted{en} ait \inserted{\unsure{un bon paquet|un tant soit peu}} est assuré par la

\uline{Proposition 9} a) Tout objet $V$ de $\mathrm{Cat}$ ayant un objet final est \unsure{tot.} $W$-asphérique -- donc \deleted{tout} \deleted{\unsure{faisceau}} \inserted{préfaisceau d'ens. $H$} \uline{représentable} sur une catégorie $U$ (la représentabilité signifiant justement que $U_{/H}$ a un objet final) est tot. $W$-asphérique.

b) \inserted{Supposons $\mathrm W(7\mathrm{bis})$ satisfait.} \deleted{Soient $V$ dans $\mathrm{Cat}$,} $V',V''$ deux ouverts de réunion $V$. Si $V',V'',V'\cap V''$ sont tot. $W$-asphériques, il en est de $=$ de $V$.

% slanted
\margin{!! Sans doute \uuline{faux} tel quel, cf rectification page suivante}

% slanted
\margin{\uuline{Faux}}

Dém \deleted{de} a). Soit $v$ \inserted{$*$} \unsure{son} objet final, et $f:V\longrightarrow V'$, \unsure{d'où} la factorisation

\diagram{\begin{tikzcd}[ampersand replacement=\&,column sep=small,row sep=small]
\& \widetilde V=\underline{\unsure{\pi|\Pi}}_0(V/\unsure{V'|U'}) \arrow[dr,"\pi"] \&\\
V \arrow[ur,"\widetilde f"] \arrow[rr,"f"'] \& \& V'
\end{tikzcd}}

de $f$. \unsure{On} \unsure{démontre|démonstration} \unsure{sans} \unsure{mal}$^{*}$ (histoire de cofinalité)

% slanted
\margin{$*$ Voici une dém. simple : pour $V\to V'$, \unsure{$=$} pour \deleted{\unsure{tout}} \deleted{\unsure{$V$}} \unsure{toute Cat} $e_V=\underset{v\ \mathrm{ds}\ V}{\varinjlim}\ v$ (\unsure{limite} inductive \unsure{du foncteur d'inclusion|des images} de $V$ dans $V^\wedge$), car $\forall x$ ds $V$, $\underset{v\ \unsure{\mathrm{ds}}\ V}{\varinjlim}\mathrm{Hom}(x,v)=\pi_0(\unsure{x\backslash V|x/V})=e$. Comme $f_!$ commute aux $\varinjlim$, on a $f_!(e_V)=\underset{v\ \mathrm{ds}\ V}{\varinjlim}f(v)$. Si $v_0$ est un obj. \unsure{final}, \unsure{cela|ceci} montre que $f_!(e_V)=f(v_0)$, \unsure{q. f.}}

\unsure{ceci} : \uline{Corollaire} : L'image par $\widetilde f$ d'un objet final de $V$ est un objet final de $\widetilde V=\underline{\pi}_0(V/U')$.

Comme ici un objet de $\mathrm{Cat}$ avec objet final est $\unsure{W}$-asphérique par $\mathrm W(3)$, cela établit a).

Prouvons b). J'ai envie de dire que si on \unsure{a} $f:V\unsure{-|\to}W$ dans $\mathrm{Cat}$, alors $\deleted{\unsure{f}}\ \underline{\pi}_0(V/w)$ \unsure{est} réunion de deux ouverts isom. à $\underline{\pi}_0(V'/w)$, $\underline{\pi}_0(V''/w)$, dont l'intersection est isom. à $\underline{\pi}_0(V'\cap V''/W)$. Cela signifierait ni plus ni
% archive sheet number: 148
